\documentclass{article}
\usepackage[utf8]{inputenc}
\usepackage[T1]{fontenc}
\usepackage{geometry}
\usepackage{graphicx}
\usepackage{float}
\usepackage{hyperref}
\usepackage[english,indonesian]{babel}

\title{Optimasi Pemrosesan Template LaTeX Menggunakan Pendekatan Berbasis AI}
\author{Academic Support Agent}
\date{\today}

\begin{document}
\maketitle

\begin{abstract}

\end{abstract}

\chapter{PENDAHULUAN}


\chapter{TINJAUAN PUSTAKA}


% Original content (commented out):
% \documentclass{article}
\usepackage[paperwidth=15cm,paperheight=15cm,margin=0.5cm]{geometry}
\usepackage[svgnames]{xcolor}
\usepackage{ifthen}

\usepackage{tikz}
\usetikzlibrary{penrose}
\usetikzlibrary{decorations.pathmorphing}
\usetikzlibrary{calc}

\colorlet{thinRhombus}{DarkOrchid}
\colorlet{thickRhombus}{DarkSlateGray}
\colorlet{circleArc}{RosyBrown}
\colorlet{longArc}{LawnGreen}

\colorlet{kite}{HotPink}
\colorlet{dart}{Fuchsia}

\colorlet{goldenTriangle}{Gold}
\colorlet{reverseGoldenTriangle}{Magenta}
\colorlet{goldenGnomon}{Cyan}
\colorlet{reverseGoldenGnomon}{LimeGreen}


\pgfmathsetmacro{\invphi}{2/(1+sqrt(5))}
\newcommand\penroseA[4]{% n, a, b, c
  \ifthenelse{0<#1}{
    {
      \pgfmathtruncatemacro{\n}{#1-1}
      \coordinate(#2#4) at ($(#2)!\invphi!(#4)$);
      \coordinate(#2#3) at ($(#2)!\invphi!(#3)$);
      \penroseA{\n}{#2}{#2#3}{#2#4}
      \penroseA{\n}{#4}{#2#3}{#3}
      \penroseB{\n}{#2#3}{#4}{#2#4}
    }
  }{
    \fill[colA,draw=white] (#2) -- (#3) -- (#4) -- cycle;
  }
}
\newcommand\penroseB[4]{% n, a, b, c
  \ifthenelse{0<#1}{
    {
      \pgfmathtruncatemacro{\n}{#1-1}
      \coordinate(#2#3) at ($(#2)!\invphi!(#3)$);
      \penroseA{\n}{#2}{#2#3}{#4}
      \penroseB{\n}{#3}{#4}{#2#3}
    }
  }{
    \fill[colB,draw=white] (#2) -- (#3) -- (#4) -- cycle;
  }
}
\begin{document}
\pagestyle{empty}
\centering

% Penrose tiling squiggles
\input{squiggly}

\clearpage 

% Penrose tiling pentagon
\null\vfill
\resizebox{\textwidth}{!}{
\begin{tikzpicture}[line width=.1pt]
  \input{pentagon}
\end{tikzpicture}
}
\vfill

\clearpage 

% Penrose Diagram
\null\vfill
\resizebox{\textwidth}{!}{
\begin{tikzpicture}
  \input{diagram}
\end{tikzpicture}
}
\vfill

\clearpage 

% Penrose Triangle
\null\vfill
\resizebox{\textwidth}{!}{
\begin{tikzpicture}
  \begin{scope}[rotate=-60, line join=bevel]
    \input{triangle}
  \end{scope}
\end{tikzpicture}
}
\vfill

\end{document}


\end{document}